\documentclass[11pt,a4paper]{article}
\usepackage[margin=2.5cm]{geometry}
\usepackage[T1]{fontenc}
\usepackage[utf8]{inputenc}
\usepackage{lmodern}
\usepackage{booktabs}
\usepackage{amsmath,amssymb}
\usepackage{microtype}
\usepackage[hidelinks]{hyperref}
\usepackage{xcolor}

\newcommand{\code}[1]{\texttt{#1}}

\title{HFP: Empirical Companion Note (v2.2)\\
\large Verified engineering claims and multi-seed recall experiments for an
$O(1)$-memory causal language-model architecture}
\author{Kayrahan Y{\i}lmaz}
\date{July 2026}

\begin{document}
\maketitle

\begin{abstract}
\noindent
This note is the empirical companion to the Hyper-Flux Projection (HFP) theory
preprint hosted in this OSF project. The theory is \emph{inspiration} for the
architecture, not evidence; conversely, nothing here validates the physics.
All claims below are machine-learning claims, tested at small scale
($\le$1M parameters, synthetic recall tasks, CPU), multi-seed where stated, and
reproducible from the public repository:
\url{https://github.com/kayra-hn/HFP} (see \code{RESULTS.md} and
\code{review\_scripts/}). Chance level is 3.3\% throughout. Effect
\emph{patterns} are the results; absolute numbers are scale-dependent.
No language-model benchmark claims are made.
\end{abstract}

\section{Architecture in one paragraph}
HFP pairs windowed local attention with a per-layer recurrent linear-attention
memory ($M \in \mathbb{R}^{D\times H}$, $z \in \mathbb{R}^{D}$) whose
inference-time state is constant in context length. Selectable axes:
retention law (\code{exp} geometric decay vs.\ \code{cubic\_flux}, an exact
discretization of $d\theta/d\tau = -\eta\,\theta^{3}$ giving state-magnitude-%
dependent decay $\lambda_t = 1/\sqrt{1+2\eta z_t^{2}}$); write rule
(\code{additive} vs.\ DeltaNet-style \code{delta}); and key feature map
(\code{elu} vs.\ DPFP, key dimension $4\times$). Engineering verified by
independent review: chunk-consistency of all mode combinations, causal
correctness, $O(1)$ state, and exactness of the cubic discretization and of its
two-pass parallel form.

\section{Methodological finding: supervision density gates learnability}
Single-query recall sequences (one supervised token per sequence) sit at the
$\ln(\mathrm{vocab})$ loss plateau and never learn once the context exceeds
$\sim$300 tokens, at any tested learning rate or curriculum---although the same
model learns the same task easily with 8 queries per sequence.
\textbf{Retention claims measured with sparse-supervision tasks are
optimization artifacts.} All experiments below use dense multi-query sequences.

\section{Retention law and write rule (3 seeds, context 160)}
Mean accuracy (\%) over seeds $\{0,1,2\}$; 600 steps; gap buckets in tokens.

\begin{table}[h]\centering
\begin{tabular}{lrrrr}
\toprule
configuration & 1--15 & 16--47 & 48--95 & 96+ \\
\midrule
exp + additive (baseline) & 44.4 & 32.7 & 18.8 & 8.1 \\
exp + delta write         & 52.1 & 31.1 & 16.4 & 6.6 \\
cubic\_flux + additive    & 31.3 & 23.0 & 11.8 & 4.7 \\
\bottomrule
\end{tabular}
\end{table}

\noindent
\code{cubic\_flux} trails the exponential baseline at this scale and is
seed-fragile (one of three seeds failed to learn); it is parked as a
long-horizon hypothesis. Delta writes help only at short range (see \S6).

\section{Length generalization (3 seeds) --- main positive result}
Models trained at context 160 are evaluated, unchanged, at $2$--$8\times$ the
training length. Fixed-gap accuracy \emph{increases} with evaluation length in
all three seeds (gap $<$48 bucket shown; \%):

\begin{table}[h]\centering
\begin{tabular}{lrrrr}
\toprule
seed & eval 160 & 320 & 640 & 1280 \\
\midrule
0 & 38.2 & 63.2 & 71.7 & 75.0 \\
1 & 32.4 & 39.3 & 45.5 & 42.9 \\
2 & 40.0 & 54.7 & 75.4 & 85.7 \\
\bottomrule
\end{tabular}
\end{table}

\noindent
The apparent ``training-length cliff'' (failure to \emph{train} at context
$\ge$320 on a fixed budget) is an optimization artifact, not an architectural
limit: the $O(1)$ recurrent state directly supports
\textbf{train-short $\rightarrow$ infer-long} deployment.

\section{The memory is interference-limited, not decay-limited}
Holding length fixed (context 640) and scaling the number of stored facts
$P = 8 \rightarrow 16 \rightarrow 24$ monotonically degrades fixed-gap accuracy
in all three seeds. This explains \S4: longer streams at fixed fact count have
lower fact density, hence less interference.

\section{Capacity axis (DPFP) --- first clear mechanism win}
DPFP (Schlag et al., 2021; key dimension $4\times$) attacks exactly the
interference limit. Confirmed across 3 seeds: at context 640, gap 256+, $P{=}8$,
the baseline is at chance in all seeds $\{5.4, 3.2, 2.5\}$ while DPFP reaches
$\{10.7, 12.9, 31.1\}$; in the highest-interference cell ($P{=}24$, gap
128--255) baseline 5.3 vs.\ DPFP 13.2. DPFP also removes the baseline's
weak-seed instability (seed-1 peak accuracy 39\% $\rightarrow$ 95\%).
It compounds with length generalization (\emph{single seed}, seed 2, trained at
160, evaluated at 1280; multi-seed confirmation of this compound cell is a
standing open item, \S8; \%):

\begin{table}[h]\centering
\begin{tabular}{lrrrr}
\toprule
variant @1280 & $<$48 & 48--127 & 128--255 & 256+ \\
\midrule
elu (baseline) & 85.7 & 69.4 & 29.9 & 5.4 \\
dpfp           & 88.1 & 87.8 & 70.1 & 33.5 \\
\bottomrule
\end{tabular}
\end{table}

\noindent
Delta writes do \emph{not} help on this task family: the interference is
cross-key feature overlap (a capacity problem), not same-key overwriting.

\section{Current recipe and honest ledger}
Recommended configuration: \code{exp} decay, \code{additive} writes,
\code{dpfp} features, standard FFN, dense multi-query training at short
context, deployment at arbitrary length. Parked or negative: \code{cubic\_flux}
(behind baseline at tested scale; exact parallel form implemented for a future
targeted long-horizon test); two-tier consolidation memory (prototype verified,
fair evaluation pending); all single-seed results are labeled as such in
\code{RESULTS.md}.

\section{Standing open items (not yet claimed)}
The following are scripted and reproducible but not yet run at the scale needed
to make a claim; none of the results above depends on them.
\begin{itemize}
\item \textbf{External family baseline.} A pure-PyTorch Gated Linear Attention
(GLA) baseline on the identical dense-retention task
(\code{review\_scripts/baseline\_compare.py}) provides an out-of-family
comparison point. Until it is run at matched parameters, all comparisons in this
note are HFP-internal (retention law / write rule / feature map), not versus the
efficient-recurrent literature.
\item \textbf{Pure-memory ablation.} Windowed attention can also see the last
\code{max\_short\_len} ring-buffer tokens in the streaming path, so ``information
travels only through $M,z$'' is verified cleanly only with the ring buffer
disabled (\code{review\_scripts/pure\_memory\_ablation.py}, ring buffer $=1$).
\item \textbf{Multi-seed DPFP$\times$length.} The compound DPFP + train-short /
infer-long cell at 1280 (\S7) is currently single-seed; seeds 0--1 pending.
\item \textbf{Language-model benchmark.} No LM-benchmark claim is made; all
results are architecture-level synthetic recall.
\end{itemize}

\section*{Reproducibility and license}
Code: \url{https://github.com/kayra-hn/HFP} (AGPL-3.0). Model card:
\url{https://huggingface.co/kayrahan35/HFP-O1-Memory-Model}. Reproduction
commands are listed in \code{RESULTS.md}; regression tests in
\code{smoke\_test.py} and \code{review\_scripts/verify\_claims.py}.

\end{document}
